\section{Introduction}

\Gls{llm} agents are increasingly employed to design complex systems, from custom silicon to controllers for fleets of autonomous robots and workflows of cooperating software agents. Agents generate hardware assertions~\cite{orenes2021autosva,kang2024fveval}, program invariants~\cite{wu2024lemur}, and code faster than design teams can review them. However, AI-generated designs are error-prone, and the errors concentrate in decomposition, which every complex design requires. In our experience, agents that decompose end-to-end hardware properties may report timed-out runs as proofs or rely on circular assumptions across lemmas. Over long design tasks, agents that delegate subtasks or compact their context restate the requirements in natural language and drift away from them~\cite{cemri2025mast}. Moreover, the reasoning behind a design is buried in transcripts, so a late error cannot be traced to its cause. As a result, AI-driven design of complex systems relies today on trial and error and on testing the final artifact. A new generation of design methods would be essential to make every step of AI-driven design traceable and checkable, and to keep the overall design correct as it is decomposed across components, agents, and abstraction levels.

Contract-based design provides a formal basis for this goal. An \gls{ag} contract specifies what a component guarantees under assumptions on its environment~\cite{benveniste2018contracts}. A decomposition is correct when the composition of the component contracts refines the system contract, and the algebra of contracts provides operators, such as the quotient, that compute the weakest specification a missing component must satisfy~\cite{incer2023algebraic,incer2022thesis}. Existing contract-based tools, including \chase~\cite{nuzzo2018chase}, Pacti~\cite{incer2023pacti}, and OCRA~\cite{cimatti2013ocra}, assume that contracts are written by human designers, and existing decomposition methods for synthesis rely on a fixed library of components~\cite{iannopollo2018decomposition}. Decomposition is also used by recent LLM reasoning methods: language models search over decompositions~\cite{yao2023tot}, compose functions validated by tests~\cite{zelikman2023parsel}, and delegate subgoals to automated provers~\cite{jiang2022thor,yang2023leandojo,ren2025deepseekprover}. However, their checks are either informal, such as tests and verbal reflection~\cite{shinn2023reflexion}, or limited to proof goals, and do not establish that the components of a system are realizable and compose correctly. Letting agents write the contracts raises three scientific questions. First, the contract operators must remain computable within the fragments that verification and synthesis tools accept; for instance, the quotient of two GR(1) contracts in general falls outside GR(1). Second, the existence of a decomposition into realizable components is undecidable for general architectures~\cite{pnueli1990distributed}, so the search over decompositions must be organized such that certified progress is never lost. Third, an agent may search indefinitely unless each failure is explained in a form that excludes whole families of candidates and improves its future proposals.

In this project, we propose a framework, termed \sysname, aiming to address these challenges by following a three-pronged approach. The central principle of \sysname is that agents propose and contracts certify: a decomposition step proposed by an agent is accepted only with a certificate that a solver can check, so that guarantees established at higher abstraction levels propagate consistently to the implementation. First, \sysname extends contract theory so that decomposition remains decidable and computable within the fragments of synthesis tools (Thrust~1). Second, it develops an agentic decomposition search in which LLM agents, guided by certified components retrieved from a library, propose contracts at every level, and only the leaf verifiers change across domains (Thrust~2). Third, it lets the agents learn from formal feedback, in context and through post-training with the verdicts of the engine as verifiable rewards (Thrust~3). \sysname builds on our prior work on contract-based design methodologies~\cite{nuzzo2015platform} and the \chase framework~\cite{nuzzo2018chase}, which provides a contract algebra with interfaces to verification and synthesis tools and to LLM agents. Our work on contract-based architecture exploration~\cite{xiao2024contrarc,xiao2026smcp} has shown that infeasibility certificates can be generalized via subgraph isomorphism to prune large design spaces, and received the Best Paper Award at DATE 2024. ACIES~\cite{xiao2026acies} applies this principle to multi-robot control, and our ongoing work to \gls{rtl} verification. If successful, this project will lay the foundations for AI-driven design of complex systems in which every decomposition proposed by an agent, for a hardware property, a robot mission, or a multi-agent workflow, is recorded as a set of contracts with a certificate that a solver can check, errors are traced to the step that caused them, and certified components are reused across designs.
